\documentclass[11pt]{amsart}

\usepackage[T1]{fontenc}
\usepackage{lmodern}
\usepackage{microtype}
\usepackage{amsmath,amssymb}
\usepackage[hidelinks]{hyperref}

\newtheorem{theorem}{Theorem}
\newtheorem{lemma}[theorem]{Lemma}
\newtheorem{corollary}[theorem]{Corollary}

\newcommand{\Zthree}{\mathbb Z/3\mathbb Z}
\newcommand{\supp}{\operatorname{supp}}

\title[Counterexamples to a packing-coloring problem]
{Two counterexamples to a packing-coloring problem of Mortada, El Zein, and Al Hajjar}

\date{}

\subjclass[2020]{05C15}
\keywords{packing coloring, claw-free graph, cubic graph, counterexample}

\begin{document}

\begin{abstract}
Mortada, El Zein, and Al Hajjar asked whether every connected claw-free
subcubic graph other than a specified graph $\mathcal H$ is
$(1,1,3,3,4)$-packing colorable. We give two nonisomorphic counterexamples,
both cubic of order $36$. Each is the triangle expansion of a cubic graph
of order $12$, and its noncolorability follows from an elementary analysis
of the possible supports of $4$-packings.
\end{abstract}

\maketitle

\section{Construction and statement}

For a positive integer $t$, a \emph{$t$-packing} in a graph is a set of
vertices at pairwise distance greater than $t$. For
$S=(s_1,\ldots,s_k)$, an \emph{$S$-packing coloring} is a partition
$V(G)=A_1\sqcup\cdots\sqcup A_k$ in which $A_i$ is an $s_i$-packing.
Mortada, El Zein, and Al Hajjar asked whether every connected claw-free
subcubic graph distinct from their graph $\mathcal H$ is
$(1,1,3,3,4)$-packing colorable~\cite[Problem~1]{MEH}.

Let $C$ be a cubic graph. Its \emph{triangle expansion} $T(C)$ is obtained
by replacing every $u\in V(C)$ with a triangle $\Delta_u$, whose three
vertices are assigned bijectively to the three edges incident with $u$,
and by joining the assigned vertices corresponding to each edge of $C$.

We define two cubic cores. The graph $C_1$ has vertex set
\[
 \{r,s,q_0,q_1,x,y,a_r,a_s,a_0,b_r,b_s,b_0\}.
\]
The set $\{q_0,q_1,x,y\}$ induces a diamond with missing edge $xy$; the
sets
\[
 A=\{a_r,a_s,a_0\},\qquad B=\{b_r,b_s,b_0\}
\]
induce triangles; and the remaining edges are
\[
 rx,\ ra_r,\ rb_r,\ sy,\ sa_s,\ sb_s,\ a_0b_0.
\]

The graph $C_2$ has vertices $p_i,q_i,x_i,y_i$ for $i\in\Zthree$. For
each $i$, the set
\[
 D_i=\{p_i,q_i,x_i,y_i\}
\]
induces a diamond with missing edge $x_i y_i$, and the remaining edges are
\[
 x_i y_{i+1}\qquad(i\in\Zthree).
\]
Put $G_1=T(C_1)$ and $G_2=T(C_2)$. For reference, graph6 encodings of
$C_1$ and $C_2$ are \verb|K?`CPagXAcAg| and \verb|K?`CPagPagEG|.

\begin{theorem}\label{thm:main}
The graphs $G_1$ and $G_2$ are nonisomorphic connected cubic claw-free
graphs of order $36$, and neither is $(1,1,3,3,4)$-packing colorable.
Consequently, Problem~1 of Mortada, El Zein, and Al Hajjar has a negative
answer.
\end{theorem}

Both cores are connected and cubic of order $12$. In a triangle expansion
the external edges form a matching, so the replacement triangles are all
the triangles. Hence $T(C)$ is connected, cubic, and claw-free, has order
$3|V(C)|$, and determines $C$ by contraction of its triangles. The core
$C_1$ has four triangles and $C_2$ has six; therefore $G_1$ and $G_2$ are
nonisomorphic. They have order $36$, whereas $\mathcal H$ has order
$12$~\cite[Figure~1]{MEH}. It remains to prove noncolorability.

\section{Distances and supports}

Index the vertices of $\Delta_u$ by the three edges incident with $u$.
The vertex indexed by an edge $e$ will be called the \emph{$e$-port} of
$\Delta_u$.

\begin{lemma}[Port criterion]\label{lem:ports}
Let $u\ne v$ be vertices of a cubic graph $C$.
\begin{enumerate}
\item If $uv\in E(C)$, then every vertex of $\Delta_u$ is at distance at
most $3$ from every vertex of $\Delta_v$ in $T(C)$.
\item If $d_C(u,v)\ge3$, then every vertex of $\Delta_u$ is at distance
at least $5$ from every vertex of $\Delta_v$.
\item Suppose $d_C(u,v)=2$. Chosen ports in $\Delta_u$ and $\Delta_v$
are at distance greater than $4$ if and only if, for every common neighbor
$w$ of $u$ and $v$, neither chosen port is indexed by the edge toward $w$.
\end{enumerate}
\end{lemma}

\begin{proof}
A simple core path $u=u_0u_1\cdots u_k=v$ lifts between chosen endpoint
ports to a path of length
\[
 2k-1+\varepsilon_u+\varepsilon_v,
\]
where each $\varepsilon$ is $0$ or $1$ according as the chosen endpoint
port is or is not indexed by the corresponding end edge of the core path.
Conversely, a shortest path in $T(C)$ that uses $m$ external edges cannot
traverse one immediately in both directions, and therefore uses at least
$m-1$ intratriangle edges. Thus its length is at least $2m-1$. The first
two assertions follow. If $d_C(u,v)=2$, every path of length at most $4$
must project to a path $uwv$ and has length
$3+\varepsilon_u+\varepsilon_v$. This gives the last assertion.
\end{proof}

For $P\subseteq V(T(C))$, define its core support by
\[
 \supp_C(P)=\{u\in V(C):P\cap\Delta_u\ne\varnothing\}.
\]

\begin{corollary}\label{cor:support}
If $P$ is a $t$-packing in $T(C)$ with $t\ge3$, then $\supp_C(P)$ is
independent in $C$. In particular, $|P|\le\alpha(C)$.
\end{corollary}

\begin{proof}
The set $P$ contains at most one vertex of each replacement triangle, and
Lemma~\ref{lem:ports}(1) excludes two adjacent core vertices from its
support.
\end{proof}

\begin{lemma}\label{lem:alpha}
One has $\alpha(C_1)=\alpha(C_2)=4$.
\end{lemma}

\begin{proof}
Let $I$ be independent in $C_1$. If $I$ contains neither $r$ nor $s$,
then it meets the diamond in at most two vertices and each of $A,B$ in at
most one. If it contains exactly one of $r,s$, say $r$, then the other
vertices lie in the three cliques
\[
 \{q_0,q_1,y\},\qquad \{a_s,a_0\},\qquad \{b_s,b_0\},
\]
so again $|I|\le4$. If it contains both $r$ and $s$, only
$q_0,q_1,a_0,b_0$ remain available, with edges $q_0q_1$ and $a_0b_0$.
Thus $|I|\le4$, and equality is attained by $\{r,s,q_0,a_0\}$.

In $C_2$, an independent set meets a diamond $D_i$ twice only in the
nonedge $\{x_i,y_i\}$. Two distinct such pairs contain one of the linking
edges $x_i y_{i+1}$, so at most one diamond is met twice. Hence
$|I|\le2+1+1=4$, with equality attained by
$\{p_0,p_1,x_2,y_2\}$.
\end{proof}

\section{The two obstructions}

Call $p_i,q_i$ the internal vertices and $x_i,y_i$ the external vertices
of the diamond $D_i$.

\begin{lemma}\label{lem:G2}
Every $4$-packing in $G_2$ has at most three vertices.
\end{lemma}

\begin{proof}
Suppose that a $4$-packing $P$ has four vertices and put
$I=\supp_{C_2}(P)$. By Corollary~\ref{cor:support} and
Lemma~\ref{lem:alpha}, $I$ is independent of size four. It therefore meets
one diamond, say $D_i$, in $\{x_i,y_i\}$ and each other diamond once.
Since $x_i$ and $y_i$ have the two common neighbors $p_i,q_i$, the port
criterion forces the selected ports in $\Delta_{x_i}$ and $\Delta_{y_i}$
to be the linking ports for $x_i y_{i+1}$ and $x_{i-1}y_i$.

Independence excludes $y_{i+1}$ and $x_{i-1}$. The two remaining external
candidates, $x_{i+1}$ and $y_{i-1}$, are adjacent. Hence one of the two
remaining support vertices is internal. If it is internal in $D_{i+1}$,
every port in its replacement triangle is at distance at most four from
the selected $x_i y_{i+1}$-port: cross to $\Delta_{y_{i+1}}$, move within
that triangle, cross to the internal vertex's triangle, and move within
it. This contradicts the $4$-packing condition. The case of an internal
vertex in $D_{i-1}$ is symmetric.
\end{proof}

Suppose that $G_2$ has a $(1,1,3,3,4)$-packing coloring with classes
$A_1,\ldots,A_5$. Each replacement triangle contains at most one vertex
of each of the independent sets $A_1,A_2$, and hence at least one vertex
of $A_3\cup A_4\cup A_5$. Therefore
\[
 |A_3|+|A_4|+|A_5|\ge12.
\]
Corollary~\ref{cor:support}, Lemma~\ref{lem:alpha}, and
Lemma~\ref{lem:G2} give the contradiction
\[
 |A_3|+|A_4|+|A_5|\le4+4+3=11.
\]

\begin{lemma}\label{lem:G1}
Let $P$ be a $4$-packing of four vertices in $G_1$, and put
$I=\supp_{C_1}(P)$. Then $C_1-I$ contains a triangle.
\end{lemma}

\begin{proof}
The support $I$ is independent. We distinguish its intersection with
$\{r,s\}$.

If $r,s\in I$, then $I$ contains one of $q_0,q_1$ and one of $a_0,b_0$.
If $a_0\in I$, the triangle $B$ lies in $C_1-I$; if $b_0\in I$, the
triangle $A$ does.

Suppose that $r\in I$ and $s\notin I$. The other three vertices of $I$
consist of one from each of
\[
 \{q_0,q_1,y\},\qquad \{a_s,a_0\},\qquad \{b_s,b_0\}.
\]
If $q_j\in I$, then the three other members of $I$ are at core distance
two from $r$ through its three distinct neighbors $x,a_r,b_r$.
Lemma~\ref{lem:ports}(3) leaves no possible port at $r$, a contradiction.
Hence $y\in I$. The pair $a_0,b_0$ is adjacent. The choice $a_s,b_0$
leaves no possible port at $a_s$, because $r,y,b_0$ are reached at core
distance two through $a_r,s,a_0$, respectively; the choice $a_0,b_s$ is
symmetric. Thus
\[
 I=\{r,y,a_s,b_s\},
\]
and $\{q_0,q_1,x\}$ is a triangle in $C_1-I$. The case
$s\in I$, $r\notin I$ is symmetric and leaves the triangle
$\{q_0,q_1,y\}$.

Finally, suppose $r,s\notin I$. Then $I$ contains $x,y$, one vertex of
$A$, and one vertex of $B$. Since $x,y$ have the two common neighbors
$q_0,q_1$, the port criterion forces their selected ports to be the ports
for $xr$ and $ys$. The chosen vertices of $A$ and $B$ must consequently
be $a_0$ and $b_0$, respectively, contradicting $a_0b_0\in E(C_1)$.
\end{proof}

Suppose that $G_1$ has a $(1,1,3,3,4)$-packing coloring with classes
$A_1,\ldots,A_5$. As above,
\[
 |A_3|+|A_4|+|A_5|\ge12.
\]
Each of these three classes has size at most four, so equality holds
throughout. Thus each has size four, and every replacement triangle
contains exactly one vertex of $A_3\cup A_4\cup A_5$. Their core supports
are therefore disjoint independent sets that partition $V(C_1)$. It
follows that
\[
 C_1-\supp_{C_1}(A_5)
\]
is bipartite, with the other two supports as a bipartition. This
contradicts Lemma~\ref{lem:G1}. Theorem~\ref{thm:main} follows.

\medskip
\noindent\textbf{Remark.}
No minimality claim is made. By Theorem~4 of~\cite{MEH}, both examples are
$(1,1,3,3,3)$-packing colorable. Thus they refute precisely the proposed
strengthening in which the last distance parameter is increased from
$3$ to $4$.

\begin{thebibliography}{9}

\bibitem{MEH}
M.~Mortada, A.~El Zein, and S.~Al Hajjar,
\newblock On $(1,1,2,3)$- and $(1,1,3,3,3)$-packing colorings of
claw-free subcubic graphs,
\newblock \href{https://arxiv.org/abs/2608.02566v1}
{arXiv:2608.02566v1} [math.CO], 3 August 2026.

\end{thebibliography}

\end{document}